% GCT_KRON_20261120, positive_formulas worker, GPT-6 Astra Pro.
% This fragment uses standard amsmath/amssymb/amsthm commands only.
% Finding means established in this run; no claim of literature originality.

\section{Positive formulas obtained from finite sets and exterior powers}
Throughout this section the ground field has characteristic zero. Write
$H_a^N=[N-a,1^a]$, $T_a^N=[N-a,a]$, and $R=(b^a)$, with $N=ab$ in the
rectangular case. Zero parts are omitted. The superscript on $H_a^N$ and
$T_a^N$ denotes the size, not a tensor power. Statements about the effective
cone of the Burnside ring are not implicit in any statement below.

\subsection{F-001: a positive kernel formula on binary contingency tables}
\textbf{Statement.}
For $0\leq r,s,u\leq N/2$, let $\mathcal T_N(r,s,u)$ be the set of eight-entry
nonnegative integral arrays $x=(x_{ijk})_{i,j,k\in\{0,1\}}$ satisfying
\[
 \sum x_{ijk}=N,\qquad \sum_{jk}x_{1jk}=r,\qquad
 \sum_{ik}x_{i1k}=s,\qquad \sum_{ij}x_{ij1}=u.
\]
Define three maps on the rational vector space with basis $e_x$ indexed by
$\mathcal T_N(r,s,u)$. The first map has target basis
$\mathcal T_N(r-1,s,u)$ and is
\[
 D_1e_x=\sum_{j,k}x_{1jk}
 e_{x-\mathbf1_{1jk}+\mathbf1_{0jk}}.                 \tag{F-001.1}
\]
A term of coefficient zero is omitted. Define $D_2,D_3$ by changing the
second and third indices in the same way. If a marginal is zero its target
space and map are zero. For the stacked matrix $D=(D_1,D_2,D_3)$,
\[
 g((N-r,r),(N-s,s),(N-u,u))
 =\dim_{\mathbb Q}\ker D
 =|\mathcal T_N(r,s,u)|-\operatorname{rank}_{\mathbb Q}D. \tag{F-001.2}
\]
In particular the coefficient is nonzero exactly when the specified integral
matrix has a nonzero rational null vector. This is a positive vector-space
formula and an exact rank criterion, not a cancellation-free enumeration of
set orbits.

\textbf{Proof.}
Let $M_v=\mathbb Q[\binom{[N]}v]$. On subset basis vectors define
$D_v(A)=\sum_{a\in A}(A\setminus\{a\})$ and
$U_v(A)=\sum_{b\notin A}(A\cup\{b\})$.
Counting a deletion followed by an insertion gives
\[
 D_{v+1}U_v-U_{v-1}D_v=(N-2v)\,\mathrm{id}_{M_v}.        \tag{F-001.3}
\]
Give subset basis vectors the standard orthonormal inner product. Then
$U_v^*=D_{v+1}$. If $v<N/2$, (F-001.3) implies
\[
 \|U_vw\|^2=\|D_vw\|^2+(N-2v)\|w\|^2,
\]
so $U_v$ is injective and $D_{v+1}$ is surjective. For
$1\leq v\leq N/2$, the characteristic of $M_v$ is $h_{N-v}h_v$ and Pieri's
rule (Macdonald, \emph{Symmetric Functions and Hall Polynomials}, second
edition, I~(5.16), p.~73, together with I~(7.3), p.~113) gives
$M_v\cong\bigoplus_{j=0}^v[N-j,j]$.
Consequently $K_v:=\ker D_v\cong[N-v,v]$; put $K_0=M_0$.
This identifies an actual irreducible kernel in a permutation module, instead
of merely subtracting two classes in its Grothendieck group.

In $M_r\otimes M_s\otimes M_u$, the intersection of the three partial
kernels is $K_r\otimes K_s\otimes K_u$. To see this directly, choose vector
space complements $M_v=K_v\oplus C_v$. The tensor product splits into eight
summands, and the first partial kernel consists of those having first factor
$K_r$, then the second and third partial kernels impose their corresponding
conditions. This argument also applies when one of the target spaces is
zero. Taking $S_N$-invariants therefore gives
\[
 \bigcap_{i=1}^3\ker(D_i|_{(M_r\otimes M_s\otimes M_u)^{S_N}})
  =(K_r\otimes K_s\otimes K_u)^{S_N}.
\]
The right side has dimension the required Kronecker coefficient.

The orbits of triples $(A,B,C)$ of subsets of prescribed cardinalities are
classified by their eight membership counts $x_{ijk}$. In the invariant
space use the \emph{average} of the orbit basis vectors,
$e_x=|O_x|^{-1}\sum_{(A,B,C)\in O_x}(A,B,C)$.
Deleting a point of type $(1,j,k)$ has $x_{1jk}$ possible choices.
Every target triple in the new orbit has the same number of preimages;
double counting gives
$|O_x|x_{1jk}=|O_{x'}|(x_{0jk}+1)$.
Thus the coefficient on orbit averages is $x_{1jk}$, exactly (F-001.1).
For unnormalized orbit sums the coefficient would instead be $x_{0jk}+1$.
This proves the formula. \hfill$\square$

\textbf{Complexity and dictionary.}
Four free entries determine the remaining entries of a binary table with
these margins, so there are at most $(N+1)^4$ tables. Rational matrix rank
yields a polynomial algorithm in $N$, with standard exact arithmetic; this
is not a claim of polynomial complexity in $\log N$.
The set of $v$-subsets is a finite constant scheme over $\mathbb Z$ carrying
an action of the constant finite group scheme $S_N$. Its free module is the
linearization of that scheme. The incidence maps are integral equivariant
maps. Over $\mathbb Q$ their kernels are representations, hence vector bundles
on the classifying stack $BS_N$, and (F-001.2) computes an invariant Hom rank
in this category. Exterior powers and Schur functors act on its $K_0$.
No claim of integral semisimplicity, integral direct-summand projectors, or a
finite-set basis for $K_v$ is made.

\subsection{F-004: a positive generating function for three two-row shapes}
\textbf{Statement.}
The following equality of formal power series holds:
\[
\begin{gathered}
 \sum_{N\geq0}\ \sum_{0\leq r,s,u\leq N/2}
 g((N-r,r),(N-s,s),(N-u,u))t^Nx^ry^sz^u\\
 =\frac{1+t^3xyz}{(1-t)(1-t^2xy)(1-t^2xz)(1-t^2yz)
 (1-t^4x^2y^2z^2)}.
\end{gathered}                                      \tag{F-004.1}
\]
The $N=0$ term is $1$. In particular, after sorting $0\leq u\leq s\leq r\leq
N/2$, put
\[
 L=\max\left(0,\left\lceil\frac{r+s+u-N}{2}\right\rceil\right),\qquad
 U=\left\lfloor\frac{s+u-r}{2}\right\rfloor.
\]
Then the coefficient is the positive integer-interval count
\[
 g((N-r,r),(N-s,s),(N-u,u))=\#([L,U]\cap\mathbb Z).       \tag{F-004.2}
\]
The interval formula is the classical result of Rosas, Corollary~5,
\emph{The Kronecker product of Schur functions indexed by two-row shapes or
hook shapes}, arXiv:math/0001084, PDF p.~11. The statement established here
is the explicit positive-series variant and its dictionary from subset
orbits; it is not a claim of originality for the coefficient formula.

\textbf{Proof of the generating variant.}
The orbit table count in F-001 has generating series
\[
 B(t,x,y,z)=\prod_{i,j,k\in\{0,1\}}(1-tx^iy^jz^k)^{-1}.
\]
For $r,s,u\leq N/2$, the coefficient of $t^Nx^ry^sz^u$ in
$(1-x)(1-y)(1-z)B$ is the Kronecker coefficient: in each factor use
$[N-r,r]=[M_r]-[M_{r-1}]$ and then take invariants. The finite differences
have additional coefficients outside this dominant chamber, so they are
not equal to (F-004.1) as an unrestricted series.

Write $G$ for the rational function on the right side of (F-004.1). For
$\epsilon=(\epsilon_1,\epsilon_2,\epsilon_3)\in\{0,1\}^3$ put
$P_\epsilon=x^{\epsilon_1}y^{\epsilon_2}z^{\epsilon_3}$, and define
\[
 \mathcal W(G)=\sum_\epsilon(-1)^{|\epsilon|}P_\epsilon
 G(tP_\epsilon,x^{1-2\epsilon_1},y^{1-2\epsilon_2},z^{1-2\epsilon_3}).
\]
There is an exact rational identity
\[
 \mathcal W(G)=(1-x)(1-y)(1-z)B.                       \tag{F-004.3}
\]
For completeness, here is a finite polynomial certificate of that identity,
which also specifies an elementary way to check every denominator and sign.
Put $K=t^2xyz$ and let $v_1=x,v_2=y,v_3=z$. Multiply (F-004.3) by
\[
 (1-K^2)\prod_\eta(1-tP_\eta)
 \prod_{i=1}^3(1-K/v_i)(1-Kv_i).
\]
All these factors are polynomials. The difference of the two sides becomes
\begin{align*}
 &\sum_\epsilon(-1)^{|\epsilon|}P_\epsilon(1+KtP_\epsilon)
  \prod_{\eta\ne\epsilon}(1-tP_\eta)
  \prod_{i=1}^3(1-Kv_i^{1-2\epsilon_i})\\
 &\hspace{10mm}-(1-x)(1-y)(1-z)(1-K^2)
   \prod_{i=1}^3(1-K/v_i)(1-Kv_i)=0.                 \tag{F-004.4}
\end{align*}
This polynomial identity follows by expanding the eight displayed summands
and collecting like monomials after $K=t^2xyz$; the supplied
\texttt{check\_formulas.py}, function \texttt{rational\_identity\_certificate},
performs exactly this multiplication in $\mathbb Z[t,x,y,z]$ and returns the
zero polynomial. Thus no convergence or numerical evaluation is involved.

Every monomial $t^Nx^ry^sz^u$ in $G$ satisfies $2r,2s,2u\leq N$: this holds
separately for every generator in its denominator and for its two numerator
monomials. A nonidentity summand of $\mathcal W$ replaces at least one such
exponent, say $r$, by $N-r+1>N/2$. Hence it contributes nothing in the
simultaneous dominant chamber. Extracting coefficients in that chamber from
(F-004.3) proves (F-004.1).

To extract its positive coefficient, choose $\varepsilon\in\{0,1\}$ from
the numerator and let $j\geq0$ be the exponent of
$t^4x^2y^2z^2$. The three exponents on pair generators are forced to be
\[
 \frac{r+s-u-2j-\varepsilon}{2},\quad
 \frac{r+u-s-2j-\varepsilon}{2},\quad
 \frac{s+u-r-2j-\varepsilon}{2}.
\]
The exponent on $t$ alone is $N-r-s-u+2j$. Integrality uniquely determines
$\varepsilon\equiv r+s+u\pmod2$. Nonnegativity of all four exponents says
exactly $L\leq j\leq U$. This proves (F-004.2). \hfill$\square$

\textbf{Geometric dictionary.}
Let $A,B,C$ be rank-two vector spaces. Schur--Weyl duality (Macdonald, I, Appendix~A, (5.5)--(5.6), p.~156) identifies the
coefficient above with the dimension of the highest-weight subspace of
$\operatorname{Sym}^N(A\otimes B\otimes C)$ of three weights
$(N-r,r),(N-s,s),(N-u,u)$. Choose rank-two free lattices with rational fibers $A,B,C$. The eight
variables are the weights of the coordinate functions on the affine dual
tensor scheme $(A\otimes B\otimes C)^*$, whose integral model is
$\mathbb A^8_{\mathbb Z}$ and whose coordinate ring is the symmetric
algebra on the tensor product of those lattices. The operator $\mathcal W$ is
rank-one Weyl antisymmetrization in each $\mathrm{GL}_2$ factor. Thus the
positive series is a multigraded highest-weight Hilbert series, while the
unrestricted table series is its torus-weight Hilbert series. This is the
precise distinction lost if one treats finite differences as positive orbit
counts.

\subsection{F-005: a symmetric positive formula for three hooks}
\textbf{Statement.}
For $0\leq r,s,u\leq N-1$,
\[
 g(H_r^N,H_s^N,H_u^N)=
 \begin{cases}
 1,&r\leq s+u,\ s\leq r+u,\ u\leq r+s,\quad r+s+u\leq2N-2,\\
 0,&\text{otherwise}.
 \end{cases}                                         \tag{F-005.1}
\]
The positive generating form is
\[
 \sum_{N\geq1}\sum_{r,s,u=0}^{N-1}g(H_r^N,H_s^N,H_u^N)t^Nx^ry^sz^u
 =\frac{t(1+t^2xyz)}{(1-t)(1-txy)(1-txz)(1-tyz)}.      \tag{F-005.2}
\]
This is a symmetric all-boundaries form of the classical hook formulas;
compare Rosas, Theorem~13(4), arXiv:math/0001084, PDF pp.~14--17. We give the
short exterior-power derivation because it exhibits exactly which
$\lambda$-operations are involved.

\textbf{Proof.}
Let $V_N=\mathbb Q^N$ be the permutation representation and
$W_N=V_N-\mathbf1$ its actual augmentation representation for $N\geq1$.
The hook identity is $\bigwedge^rW_N\cong H_r^N$: indeed,
$\bigwedge^rV_N$ is induced from $\operatorname{sgn}_{S_r}\boxtimes\mathbf1$,
so the vertical Pieri formula gives $H_r^N\oplus H_{r-1}^N$ in the
interior degrees; induction using $V_N=\mathbf1\oplus W_N$ gives the
claim, with the endpoints checked separately. The exact Pieri locator is
Macdonald, I~(5.17), p.~73.
A permutation with cycle lengths $m_1,\ldots,m_q$ satisfies
\[
 \lambda_x(V_N)(\sigma)=\det(1+x\sigma|V_N)
 =\prod_{i=1}^q(1-(-x)^{m_i}).
\]
For each cycle type, averaging the product of the three such determinants
computes the invariant dimension. The symmetric-group cycle-index identity
therefore gives
\begin{align*}
 R(t,x,y,z)
 &=\sum_{N\geq0}t^N\frac1{N!}\sum_{\sigma\in S_N}
   \det(1+x\sigma)\det(1+y\sigma)\det(1+z\sigma)\\
 &=\exp\left(\sum_{m\geq1}\frac{t^m}{m}
 (1-(-x)^m)(1-(-y)^m)(1-(-z)^m)\right)\\
 &=\frac{(1+tx)(1+ty)(1+tz)(1+txyz)}
 {(1-t)(1-txy)(1-txz)(1-tyz)}.
\end{align*}
The exponential identity follows directly by choosing the number of cycles
of each length; a cycle type with $c_m$ cycles of length $m$ has normalized
weight $\prod_m(m^{c_m}c_m!)^{-1}$. Thus it also proves the only cycle-index
fact used here.

Removing the trivial summand in each exterior-power series divides by
$(1+x)(1+y)(1+z)$ for $N\geq1$. The $N=0$ term must first be removed.
The elementary factorization
\begin{align*}
 &(1+tx)(1+ty)(1+tz)(1+txyz)
 -(1-t)(1-txy)(1-txz)(1-tyz)\\
 &\hspace{15mm}=t(1+x)(1+y)(1+z)(1+t^2xyz)
\end{align*}
proves (F-005.2).

For coefficient extraction, choose $\varepsilon=0$ or $1$ in the numerator.
The pair exponents are forced to be
$(r+s-u-\varepsilon)/2$, $(r+u-s-\varepsilon)/2$,
$(s+u-r-\varepsilon)/2$; their integrality uniquely determines
$\varepsilon\equiv r+s+u\pmod2$.
They are nonnegative exactly when the triangle inequalities hold. The
remaining exponent of $t$ is
$N-1-(r+s+u+\varepsilon)/2$; it is nonnegative exactly when
$r+s+u\leq2N-2$. Each admissible coefficient is therefore counted once.
The formula also shows $r,s,u\leq N-1$ automatically, and covers $N=1$.
\hfill$\square$

\textbf{Dictionary.}
Over $\mathbb Z$, $W_N$ is the locally free kernel of the augmentation map
$\mathbb Z^N\to\mathbb Z$. It is a representation of the constant group
scheme $S_N$, and $\lambda^r[W_N]=[\bigwedge^rW_N]$ in the representation
category's $K_0$. Over $\mathbb Q$ the invariant projection is the Reynolds
idempotent $N!^{-1}\sum\sigma$. Formula (F-005.2) computes ranks of these
invariant vector bundles. These are genuine exterior powers, with the
orientation signs of permuted wedges. They must not be replaced by the
unsigned set of $r$-subsets.

\subsection{F-002: the positive stable rectangle--two-row formula}
\textbf{Statement and source.}
Let $p_r(a,b)$ count partitions of $r$ in an $a$ by $b$ rectangle, and put
$p_{-1}=0$. The known rectangle identity is
\[
 g((b^a),(b^a),(ab-r,r))=p_r(a,b)-p_{r-1}(a,b),
 \qquad 0\leq r\leq ab/2.                            \tag{F-002.1}
\]
Its exact source is Pak--Panova, \emph{Strict unimodality of $q$-binomial
coefficients}, Lemma~3, p.~2. Suppose $a\leq b$. In the larger stable window
$r\leq b$ (still with $r\leq ab/2$), the positive refinement is
\[
 g((b^a),(b^a),(ab-r,r))
 =[q^r]\prod_{i=2}^a(1-q^i)^{-1}.                    \tag{F-002.2}
\]
Thus it counts partitions of $r$ into parts from $\{2,3,\ldots,a\}$.
For $a=1$ it is $1$ only at $r=0$; for $a=2$ it is $1$ precisely for even
$r$; for $a\geq3$ it is positive precisely for $r=0$ or $r\geq2$ in this
window.

\textbf{Proof of the refinement and the Burnside dictionary.}
Since $r\leq b$, the width condition in the rectangle imposes no restriction.
Partitions counted by $p_r(a,b)$ have at most $a$ parts. Conjugation
identifies them with partitions whose parts are at most $a$.
Multiplication of their generating series by $(1-q)$ removes the arbitrary
number of parts equal to $1$, proving (F-002.2). Its support statements
follow from the empty partition, parity for $a=2$, and the fact that every
integer $r\geq2$ is $2v+3w$ with $v,w\geq0$.
For the link with F-001, the two-row identity
$[N-r,r]=[M_r]-[M_{r-1}]$ and Frobenius reciprocity give
\[
 g(\alpha,\beta,(N-r,r))=A_r(\alpha,\beta)-A_{r-1}(\alpha,\beta),\qquad
 A_r=\sum_{\eta\vdash r,\theta\vdash N-r}
 c^\alpha_{\eta\theta}c^\beta_{\eta\theta}.
\]
For a rectangle $R=(b^a)$, the coefficient $c^R_{\eta\theta}$ is $1$ exactly
when $\theta$ is the rotated complement of $\eta$ in $R$, and is $0$
otherwise. One way to check this particular Littlewood--Richardson fact is
$S_R(\mathbb Q^a)=\det^b$: its multiplicity in
$S_\eta\otimes S_\theta$ is the multiplicity of $S_\theta$ in
$\det^b\otimes S_\eta^*$, an irreducible whose weight is that rotated
complement if it is polynomial. Thus $A_r(R,R)=p_r(a,b)$.
This recovers the known identity in precisely the form needed for the
finite-set linearization. \hfill$\square$

\subsection{F-003: rectangular hooks, with all endpoints}
\textbf{Statement and source.}
Put $m=\min(a,b)$. For every $0\leq k\leq ab-1$,
\[
 g((b^a),(b^a),(ab-k,1^k))
 =[q^k]\prod_{i=2}^m(1+q^{2i-1}).                    \tag{F-003.1}
\]
This is Ikenmeyer--Panova, \emph{Rectangular Kronecker coefficients and
plethysms in geometric complexity theory}, Proposition~6.4, PDF p.~15;
attached source lines 1380--1427. We use the empty-product convention:
if $m=1$, only $k=0$ has value $1$. The polynomial has degree $m^2-1$, so
all $k>m^2-1$ vanish, even when the rectangle is much longer in one direction.
The nonvanishing criterion is exact: $k$ must be a sum of distinct members
of $\{3,5,\ldots,2m-1\}$.

\textbf{Short exterior-power verification of the boundaries.}
The permutation representation satisfies
$\bigwedge^kV_N=H_k^N\oplus H_{k-1}^N$ for $1\leq k\leq N-1$,
with the obvious single summands at $k=0,N$ and $H_{-1}=H_N=0$.
The module $\bigwedge^kV_N$ is induced from the sign representation on
$S_k$ and the trivial representation on $S_{N-k}$, not the unsigned subset
module. Restricting the two rectangular factors and using the complement
rule from F-002 gives
\[
 g_k+g_{k-1}=\#\{\eta\vdash k:\eta=\eta',\ \eta\subseteq(b^a)\}.
\]
Self-conjugate diagrams inside the rectangle are exactly those inside the
$m$ by $m$ square. Their diagonal hook lengths are distinct odd integers in
$\{1,3,\ldots,2m-1\}$, and these lengths sum to the diagram's size.
This bijection gives $(1+q)\sum g_kq^k=\prod_{i=1}^m(1+q^{2i-1})$.
Cancelling $1+q$ proves the endpoint extension and records explicitly the
orientation needed in the $\lambda$-operation. \hfill$\square$

\subsection{F-006 and F-007: a corrected hook-positivity statement}
\textbf{F-006 (proved refinement).}
For $m\geq7$, the support of the product in (F-003.1) is
\[
 \{0,1,\ldots,m^2-1\}\setminus
 \{1,2,4,6,m^2-7,m^2-5,m^2-3,m^2-2\}.               \tag{F-006.1}
\]
Here is an elementary full proof of the support assertion, independent of
any strict-unimodality assertion. Let $A_m=\{3,5,\ldots,2m-1\}$ and
$D_m=\sum A_m=m^2-1$. Subset complementation makes its subset-sum support
symmetric about $D_m/2$. No subset sums to $1,2,4,6$; their complements are
also impossible. For $m=7$, all values $7$ through $24$ are represented by
\[
\begin{array}{c|rrrrrrrrr}
k&7&8&9&10&11&12&13&14&15\\ \hline
\text{sum}&7&3+5&9&3+7&11&3+9&13&3+11&3+5+7
\end{array}
\]
\[
\setlength{\arraycolsep}{3pt}
\begin{array}{c|rrrrrrrrr}
k&16&17&18&19&20&21&22&23&24\\ \hline
\text{sum}&3+13&3+5+9&5+13&3+5+11&7+13&3+5+13&9+13&3+7+13&11+13.
\end{array}
\]
Their complements cover $24$ through $41=D_7-7$. The additional values
$0,3,5,D_7-5,D_7-3,D_7$ are immediate. This proves the base case.
If the assertion holds for $m\geq7$, the new part $2m+1$ shows that the next
support contains both intervals
$[7,D_m-7]$ and $[2m+8,D_{m+1}-7]$. They overlap or are adjacent because
$2m+8\leq m^2-7$ for $m\geq7$. The boundary values $0,3,5$ and their
new complements are present. The eight impossible values already identified
are the only remaining values. This proves (F-006.1). \hfill$\square$

\textbf{F-007 (refutation of literal source wording).}
In the attached version of Ikenmeyer--Panova Theorem~6.1
(\texttt{thm:hook\_positive}, source lines 1340--1356), the last exceptional
value is printed as $d^1-2$. The correct last value is $d^2-2$, as follows
both from Proposition~6.4 and from the same source's later table.
More substantively, its unqualified sentence
$g_k(d,n)>g_{k-2}(d,n)$ for $k\leq d^2/2$ is false. For every $d\geq2$ and
$n\geq d$, (F-003.1) gives $g_0=1$ and $g_2=0$; for $d\geq2$, $k=2$ is
in the stated range. In particular the failure persists for $d\geq7$.
This refutes that sentence, not Proposition~6.4 or the corrected support
classification. We use F-006, not the false strictness claim, anywhere that
hook support is needed.

\subsection{F-008: the Grassmannian interpretation and its limit}
Let $X=\operatorname{Gr}(a,a+b)$ over $\mathbb Z$, of dimension $ab$.
Its Schubert cells give
$\dim H^{2r}(X_{\mathbb C},\mathbb Q)=p_r(a,b)$ and vanishing odd cohomology.
Cup product by the Pl\"ucker hyperplane class $L$ is injective from
$H^{2r-2}$ to $H^{2r}$ for $r\leq ab/2$, by hard Lefschetz. Hence
\[
 g((b^a),(b^a),(ab-r,r))
 =\dim\frac{H^{2r}(X_{\mathbb C},\mathbb Q)}{L H^{2r-2}(X_{\mathbb C},\mathbb Q)}.
                                                            \tag{F-008.1}
\]
The quotient is isomorphic to primitive cohomology after choosing the
Lefschetz decomposition. Formula (F-008.1) is a consequence of the known
rectangle identity and hard Lefschetz; it asserts equality of dimensions,
not a constructed canonical identification with a Kronecker multiplicity
space. Over a finite field the Schubert cell filtration gives the same
ranks in $\ell$-adic cohomology, and the even groups have Tate Frobenius
weights $2r$; the normalized quotient after a Tate twist has weight zero.
These facts provide the requested schemes, cohomology, weights, and Tate
motive dictionary in this family. They do not prove strict positivity of
every primitive degree or a saturation theorem: the parity example with
$a=2$ is an explicit obstruction. None of this argument uses RH or GRH.
